\documentclass[11pt]{article}

% =============================================================
%  The Reliability Tax — ACL 2026 course project (Research track)
%  Draft: Abstract / Introduction / Related Work
%
%  NOTE ON PLACEHOLDERS: experiments have NOT been run yet.
%  Every empirical claim is marked with \PH{...} and must be
%  replaced with measured values before submission. Do not
%  fabricate numbers.
% =============================================================

\usepackage[review]{acl}
\usepackage{times}
\usepackage{latexsym}
\usepackage[T1]{fontenc}
\usepackage[utf8]{inputenc}
\usepackage{microtype}
\usepackage{inconsolata}
\usepackage{graphicx}
\usepackage{amsmath,amssymb}

% placeholder macro for un-run results (delete before camera-ready)
\newcommand{\PH}[1]{\textbf{[\,#1\,]}}

\title{The Reliability Tax: Quantifying the Capability--Reliability Trade-off in Reasoning LLMs}

\author{Anonymous ACL submission}

\begin{document}
\maketitle

\begin{abstract}
Reasoning-augmented language models (LRMs) are now deployed by default for
high-stakes queries, and test-time compute---how long a model is allowed to
``think''---is increasingly a user-facing knob.
The prevailing assumption is that more thinking yields more trustworthy
answers, but recent work has begun to undermine it: sweeping the reasoning
budget impairs confidence calibration, and the marginal utility of extra
reasoning turns negative.
These findings are single-axis and descriptive: they show that
the trade-off \emph{exists}, not whether it is uniform or---the property a
deployment rule needs---whether it is \emph{predictable}.
We ask whether a black-box user can foresee, before committing to a budget,
that more thinking will hurt.
Our budget-sweep protocol varies thinking across six models from four
families over four task types, tracking accuracy and three reliability axes:
calibration, selective prediction, and
abstention versus fabrication on false-premise items, where \emph{fabrication}
means endorsing a premise that is in fact false.
We first measure the trade-off and find it is \emph{task-dependent}:
\PH{the ratio of reliability loss to accuracy gain varies by} $\PH{x}\times$
across tasks.
We then show it is \emph{partly predictable}: using only the trace and
verbalized confidence observed at a low budget, a lightweight predictor
decides, on held-out items, whether continuing will help; the resulting
oracle-free stopping rule attains higher \emph{usable accuracy}---accuracy at
a fixed reliability floor---than the best fixed budget at matched cost, while
a permutation control on the trace features loses the advantage.
These results recast ``think longer'' from a default into a decision.
\end{abstract}

\section{Introduction}

Reasoning-augmented language models (LRMs) such as o1-style systems and
DeepSeek-R1 have become the default for tasks that were previously out of
reach, and their defining property is that they spend a variable amount of
compute \emph{thinking before answering} \citep{Guo2025DeepSeekR1}.
This has turned test-time compute from an implementation detail into a
first-class, increasingly user-facing control: models expose reasoning
budgets, ``thinking effort'' levels, and explicit toggles
\citep{Muennighoff2025s1, Snell2024TestTimeCompute}.
The implicit contract behind these knobs is simple and appealing: \emph{more
thinking should buy more trustworthy answers}.
Practitioners act on it---escalating to higher reasoning budgets, longer
chains of thought, or a reasoning model over a chat model whenever an answer
matters.
A customer-support or clinical-QA pipeline, for instance, may simply raise the
reasoning budget to ``max'' for every incoming query, on the assumption that
the most careful answer is the safest.
The cost of that assumption is concrete: an over-confident wrong answer is
worse than an admitted uncertainty, because nothing downstream flags it, and a
model that fabricates support for a false premise is more dangerous than one
that declines to answer.

This contract is, however, being contradicted from several directions at
once.
Training that raises reasoning accuracy with verifiable rewards
systematically degrades \emph{calibration}: incorrect answers acquire
\emph{lower} perplexity than correct ones, so a model's confidence stops
tracking its correctness \citep{Wang2026CAPO, Jin2025Entropy}.
Enhancing reasoning amplifies tool hallucination---models fabricate
non-existent tools---and the effect survives even when the enhancement is
purely inference-time \citep{Yin2025ReasoningTrap}.
In multimodal settings, ``slower'' models are measurably \emph{less} truthful
under misleading inputs \citep{Fang2025Slower}.
Reasoning models commit to false premises more decisively than their
non-reasoning counterparts, becoming \emph{confidently wrong} on 26--42\% of
false-premise queries \citep{Sathyanathan2025Presuppositions}, and even
answers that are perfectly self-consistent can collapse under mild
contextual interference \citep{Xu2026IllusionConfidence}.
Most directly, sweeping the reasoning budget itself---at inference time, with
no training involved---has been observed, in several settings, to impair
confidence calibration \citep{DontThinkTwice2025} and to yield non-monotonic
marginal returns \citep{WhenMoreThinkingHurts2026}.

Each of these is a warning; none is a law.
The existing evidence is fragmented along three axes.
First, it is \emph{task- and family-local}: a single benchmark, often a
single model or training recipe, is used to establish the effect.
Second, it is overwhelmingly \emph{training-side}: the sharpest results
concern what RLVR does to weights \citep{Wang2026CAPO, Jin2025Entropy},
which is unavailable to the many practitioners who only have API access.
Third, reliability is measured \emph{one axis at a time}---calibration
\emph{or} selective prediction \emph{or} abstention---so it is impossible to
tell whether these failures share a cause or merely co-occur.
More fundamentally, even the closest work is \emph{descriptive}: none reports
whether the three axes degrade \emph{together}, as one underlying cost, or
whether the degradation is \emph{predictable} from the reasoning trace before
a budget is committed---the property a deployment rule actually needs.
Consequently, we cannot currently answer the question that actually governs
deployment: \emph{does increasing a model's thinking budget systematically
trade reliability for accuracy, and is that trade-off universal or
task-specific?}

We address this gap with a controlled, black-box study.
We formalize the trade-off as a \emph{reliability tax}: with the thinking
budget $b$ as the controlled variable and every reliability score $\rho_r$
oriented so that higher is better, we estimate the ratio of reliability change
to accuracy change,
$\tau_r = (\partial \rho_r / \partial b) / (\partial \mathrm{acc} / \partial b)$,
for each reliability axis $r$; a large $\tau_r$ means accuracy is bought at a
steep reliability price---exactly the quantity a deployment should minimize.
Because $\tau_r$ diverges where accuracy saturates, our primary summaries are
bounded instead: the fixed-coverage risk and the area under the
risk--coverage curve, with $\tau_r$ reported only where
$\partial \mathrm{acc}/\partial b > 0$.
The three axes correspond to the three decisions a high-stakes deployment must
make---\emph{how much to trust an answer} (calibration), \emph{how much
coverage to serve} (selective prediction), and \emph{whether to answer at all
when the premise is false} (abstention versus fabrication)---and we measure
them jointly precisely to test whether they degrade as one or independently.
Our \emph{budget-sweep protocol} holds the prompt set fixed and varies a
model's thinking budget across its supported range (reasoning-effort levels,
reasoning-token caps, and black-box budget forcing \citep{Muennighoff2025s1}),
tracking, at each budget, accuracy and the three axes: calibration, via ECE
and verbalized overconfidence \citep{Kadavath2022KnowWhatTheyKnow,
Tian2023JustAsk}; selective prediction, via the risk--coverage curve; and the
\emph{fabricate-versus-abstain} choice on false-premise items
\citep{Sathyanathan2025Presuppositions, Xu2026Presupposition}.
We sweep six models from four families---open-weight and closed API, all
exposing a thinking toggle or a reasoning-effort level---over four task types
chosen to span where extra compute should help and where it can hurt:
mathematical reasoning (verifiable), factual QA (closed-book knowledge),
multi-hop QA (composition), and false-premise items (where the model should
refuse).
Because budget knobs are not comparable across providers, we report against
the \emph{realized} number of reasoning tokens and include an equal-compute
control; each item is rendered under several prompt templates so that
form-induced variance can be separated from sampling variance
\citep{Miller2024ErrorBars, Zhang2025BeyondSingular}; and all calls are pinned
to dated model snapshots, since providers update models silently.
Where a provider does not expose raw reasoning traces, our predictor falls
back to the answer and verbalized confidence alone, so the protocol never
assumes white-box access.
We define \emph{usable accuracy} as the accuracy attainable at a fixed
reliability floor---the accuracy over the subset served at the coverage where
calibration error stays below a target $\epsilon$---so that a budget is judged
by what it can be trusted to serve, not by its best-case answer.
This design tests one descriptive and two predictive claims:
\emph{H1 (task-dependence)} the ratio of reliability loss to accuracy gain
differs across task types by at least a factor of two---falsified if the
ratios are statistically indistinguishable;
\emph{H2 (predictability)} a predictor trained on a held-out half decides, at
low budget and without seeing the high-budget answer, whether continuing will
help, and the induced policy beats the best fixed budget in usable accuracy at
matched cost---falsified otherwise;
\emph{H3 (negative control)} shuffling the trace features removes H2's
advantage---if the shuffled control matches, the predictor exploits an
artifact rather than the trace.

Our contributions are:
\begin{itemize}
  \item A black-box \emph{budget-sweep protocol} and a multi-axis reliability
  measurement---calibration, selective prediction, and
  abstention-versus-fabrication---that any API user can apply without weights
  or training.
  \item To our knowledge, the first cross-task and cross-family test of
  whether the trade-off is a law or a set of local effects, and of where it
  is worst.
  \item A held-out, cost-matched evaluation of an \emph{oracle-free stopping
  rule} that decides when to stop thinking, together with a permutation
  control that rules out trace-format artifacts.
  \item The protocol, prompts, and raw responses, plus a reproducibility
  recipe that pins dated API snapshots.
\end{itemize}

\section{Related Work}

\paragraph{Reasoning models and test-time compute.}
LRMs are trained to emit long reasoning traces before answering, most
prominently through reinforcement learning with verifiable rewards
\citep{Shao2024GRPO, Guo2025DeepSeekR1}.
A parallel line studies how to spend \emph{inference-time} compute: scaling
test-time search against verifiers can outperform a larger base model
\citep{Snell2024TestTimeCompute}, and simple ``budget forcing'' can extend or
truncate a trace to trade quality for cost \citep{Muennighoff2025s1}.
A complementary strand asks what these traces actually buy.
\citet{Shojaee2025IllusionThinking} show accuracy collapses past a complexity
threshold and that reasoning effort rises then \emph{falls} as problems get
harder; \citet{Srivastava2025Overthink} find reasoning models generate roughly
$18\times$ more tokens while sometimes being \emph{less} accurate, with
non-monotonic returns to budget.
These works motivate treating the thinking budget as the independent variable,
but they evaluate final-answer accuracy rather than reliability, and recent
work further questions whether RLVR expands capability at all
\citep{Yue2025RLBeyondBase, Zhao2025SpuriousRewards}.

\paragraph{Efficient and adaptive reasoning.}
A large literature instead makes reasoning \emph{cheaper}, by truncating or
compressing traces and by early exit \citep{Chen2024Overthinking,
Sui2025StopOverthinking, StepGRPO2026, EarlyExit2026}, while a recent strand
characterizes the marginal utility of extra tokens
\citep{WhenMoreThinkingHurts2026}.
Early-exit methods share our goal of deciding \emph{when to stop}, but differ
in both setting and signal: they read the model's \emph{internal} hidden states
and are therefore white-box, and they optimize cost subject to accuracy.
We stop from the \emph{external} trace and verbalized confidence alone---the
only signals an API user has---and we optimize \emph{trust}: the question is
not whether the answer stays as accurate, but whether it stays as reliable.

\paragraph{Calibration, selective prediction, and abstention.}
A well-calibrated model's confidence predicts its correctness, which is what
makes deferral and abstention meaningful.
Early work showed that conditional probabilities of large pretrained models
are reasonably calibrated and that models can verbalize calibrated confidence
\citep{Kadavath2022KnowWhatTheyKnow, Lin2022Verbalized}, while
later work finds that instruction tuning and RLHF harm calibration and that
\emph{verbalized} confidence is often the better estimate
\citep{Tian2023JustAsk, Xiong2023Confidence, SanzGuerrero2026Overconfident}.
Most recently, training-side analyses show RLVR makes models overconfident,
degrading relative calibration as accuracy rises
\citep{Wang2026CAPO, Jin2025Entropy}, and that superficially consistent
beliefs can be brittle \citep{Xu2026IllusionConfidence}.
The closely related notion of \emph{abstention}---knowing when not to
answer---dates back to selective classification \citep{Geifman2017Selective}
and is surveyed for LLMs by \citet{Slobodkin2024Abstention}; strikingly,
reasoning models are \emph{worse} at abstaining on unanswerable questions
\citep{Kirichenko2025AbstentionBench}.
Calibration itself is usually operationalized through ECE, introduced for
neural networks by \citet{Guo2017ECE}.
This literature establishes \emph{that} calibration can degrade with
reasoning, but studies one family, one training recipe, or one axis; whether
calibration, abstention, and fabrication move together as a function of the
budget a black-box user selects is unexamined.

\paragraph{Faithfulness, sycophancy, and false premises.}
Reasoning traces are not always faithful to the computation that produced
the answer \citep{Turpin2023Unfaithful, Chen2025DontSay}, and
models are readily steered by user framing or social pressure
\citep{Sharma2023Sycophancy}.
Hallucination is a growing concern: models fabricate tools as reasoning
strengthens \citep{Yin2025ReasoningTrap}, produce unsupported claims that
snowball across steps \citep{Zhang2023Snowball}, and cannot reliably
self-correct their own reasoning \citep{Huang2024SelfCorrect}.
On false-premise queries, reasoning models are confidently wrong on a large
fraction of items \citep{Sathyanathan2025Presuppositions}, and do not reliably
reason about the presuppositions such queries turn on
\citep{Xu2026Presupposition}; multimodal reasoning models become less truthful
as they think longer \citep{Fang2025Slower}.
Each is typically studied in isolation; none is probed as a function of the
thinking budget.

\paragraph{Reliability pitfalls in evaluation.}
Measuring reliability is itself error-prone.
Benchmark scores carry large sampling variance---on GSM8K a single generation
can vary by 18.6 points across runs---so multi-sample, hierarchical
estimation is needed \citep{Zhang2025BeyondSingular, Miller2024ErrorBars}.
We therefore adopt multi-sample estimation throughout.

\paragraph{Closest work and our position.}
The closest studies vary the reasoning budget and measure a single
reliability axis: \citet{DontThinkTwice2025} show that over-reasoning
impairs confidence calibration, and \citet{WhenMoreThinkingHurts2026}
characterize the diminishing, sometimes negative, marginal utility of extra
tokens.
Both are single-axis, and neither asks whether the losses across axes
\emph{co-move} or can be \emph{predicted}.
We therefore change the question from \emph{whether} the trade-off exists to
\emph{whether it is a law}: we measure calibration, selective prediction, and
fabrication jointly across tasks and families, and test whether
impending degradation is foreseeable from the trace, turning the trade-off
curve into a black-box budget-selection rule.

\section{The Reliability Tax}

Everything in this paper is read off a single object: the
\emph{budget--reliability profile} of a model on a task---how accuracy and
each reliability axis change as the model is allowed to think longer.
We fix notation and define this object (\S\ref{sec:profile}), read the tax
from it (\S\ref{sec:tax}), operationalize the three axes
(\S\ref{sec:axes}), and turn the profile into a decision rule
(\S\ref{sec:rule}).

\subsection{The budget--reliability profile}\label{sec:profile}
Let $M$ be a model and $\mathcal{B}=\{b_0<\dots<b_B\}$ its supported thinking
budgets---reasoning-effort levels, reasoning-token caps, or black-box budget
forcing \citep{Muennighoff2025s1}.
For a prompt $x$ and budget $b$, the model returns a response $y\sim M(x,b)$
carrying a \emph{realized} reasoning length $\ell(y)$ in tokens.
Because a nominal budget means different things to different providers, we plot
every quantity against $\ell$ rather than $b$.
For a fixed prompt set $\mathcal{D}$, the profile is the collection
\begin{equation}
\big\{\mathrm{acc}(b),\ \rho_1(b),\ \rho_2(b),\ \rho_3(b)\big\}_{b\in\mathcal{B}},
\end{equation}
estimated with $k$ samples per $(x,b)$ and reported with confidence intervals
clustered by prompt---prompts, not samples, are the unit of variation
\citep{Zhang2025BeyondSingular, Miller2024ErrorBars}.
A profile is thus a two-dimensional object, budget against a reliability
reading, and the paper is a study of its shape.

\subsection{The reliability tax}\label{sec:tax}
Orient each reliability score $\rho_r$ so that higher is better, and define the
tax along the profile as
\begin{equation}
\tau_r(b)=\frac{\partial \rho_r/\partial b}{\partial \mathrm{acc}/\partial b},
\end{equation}
the reliability given up per unit of accuracy gained.
A deployment wants $\tau_r$ small; a large $\tau_r$ means extra accuracy is
bought at a steep reliability price.
$\tau_r$ is undefined where accuracy saturates---precisely the regime we care
about---so our primary summaries are \emph{bounded}: the fixed-coverage risk
$R(c)$, the error rate over the top-$c$ fraction of the most confident
answers, and the area under the risk--coverage curve \citep{Geifman2017Selective};
we report $\tau_r$ only where $\partial\mathrm{acc}/\partial b>0$.

\subsection{Three reliability axes}\label{sec:axes}
The axes are the three decisions a high-stakes deployment must make.
\paragraph{Calibration (how much to trust an answer).}
Expected calibration error \citep{Guo2017ECE} and \emph{overconfidence}, the
mean verbalized confidence on incorrect answers
\citep{Kadavath2022KnowWhatTheyKnow, Tian2023JustAsk, Xiong2023Confidence}.
\paragraph{Selective prediction (how much coverage to serve).}
$R(c)$ and AURC, which bound how much of the task a system can answer before
its error rate exceeds a target.
\paragraph{Abstention versus fabrication (whether to answer a false premise).}
On false-premise items we classify each response as \emph{refute},
\emph{hedge}, or \emph{endorse}, and report the fabrication rate---endorsing a
premise that is false \citep{Sathyanathan2025Presuppositions,
Kirichenko2025AbstentionBench}.
The three are measured jointly because the question is not whether any one
degrades, but whether they degrade \emph{together}.

\subsection{From profile to decision: a stopping rule}\label{sec:rule}
The profile also yields an oracle: $b^\star$ maximizes usable accuracy, the
accuracy at the coverage where calibration error stays below a target
$\epsilon$.
The actionable question is whether $b^\star$ can be found \emph{without}
running the expensive end of the sweep.
We train a predictor on a held-out half of the prompts: at a low budget $b_0$
it observes features $f(y)$---realized length, verbalized confidence,
self-consistency across $k$ samples, and hedging markers---and decides whether
continuing to the top budget will help.
The induced policy is evaluated on the other half, compared against every
fixed budget under matched cost, and stress-tested by shuffling $f$ across
prompts: if the shuffled control matches, the predictor exploits an artifact
rather than the trace.
Where a provider does not expose raw reasoning traces, $f$ reduces to the
answer and its verbalized confidence, so the rule never assumes white-box
access.

\section{Experimental Setup}
\paragraph{Models.}
Six models from four widely-deployed families, spanning open-weight and
closed API, each exposing a thinking toggle or a distinct reasoning variant:
DeepSeek (R1 and V3), Qwen3 (235B and 32B, with an \texttt{enable\_thinking}
toggle), GLM-4.6, and Kimi-K2.
All are reachable through OpenAI-compatible endpoints, and each provides at
least two settings of the budget $b$, so that $b$ is a real knob rather than a
prompt artefact.
\paragraph{Tasks.}
Four types chosen to span where extra compute should help and where it can
hurt: mathematical reasoning (GSM8K, MATH-500; verifiable), factual QA
(TriviaQA, NQ; closed-book knowledge), multi-hop QA (HotpotQA; composition),
and \emph{false-premise} items, drawn from fact-checked claims known to be
false \citep{Wadden2020SciFact} and rewritten as questions that presuppose
them (e.g., ``Why did Einstein win the 1921 Nobel Prize in Literature?''),
where the reliable behaviour is to refute rather than answer. \PH{200 items
per task.}
\paragraph{Budgets and sampling.}
\PH{Four budget levels} per model, \PH{$k=4$} samples per cell at temperature
$0.7$, under \PH{two} prompt templates so that form-induced variance can be
separated from sampling variance \citep{Miller2024ErrorBars}.
\paragraph{Statistics.}
Prompt-clustered bootstrap intervals and mixed-effects models with prompt and
model as random effects; the negative control is a permutation test over $f$.
\paragraph{Reproducibility.}
All calls are pinned to dated model snapshots; prompts, responses, and scripts
are released.

\subsection{A minimal viable pilot}
Before the full sweep we run the smallest end-to-end version that exercises
every stage---budget sweep, profile, tax, stopping decision---on a single
model and two budgets (\PH{100 items}, \PH{$k=4$}).
To validate the pipeline itself rather than a model, we first drive it with a
synthetic backend that plants a known tax (accuracy rises modestly while
calibration and fabrication degrade); the pipeline recovers the planted signs
and selects the lower budget as the trustworthy one, and the same script runs
unchanged against a real API.
The pilot is the unit of validation: if the pipeline cannot recover a known
effect here, the larger sweep will not either.

\section{Results and Analysis}
\paragraph{RQ1: Is the tax real, and is it a law?}
\PH{Profile shapes per model and task; tax magnitude; and whether the ratio of
reliability loss to accuracy gain differs across tasks by the factor of two
predicted by H1.}
\paragraph{RQ2: Is it predictable?}
\PH{Predictor accuracy on held-out prompts; the stopping rule's usable
accuracy against the best fixed budget under matched cost (H2); and the
shuffled-feature control (H3).}
\paragraph{RQ3: What does it cost in practice?}
\PH{The share of queries on which the default of ``max budget'' over-thinks,
and the usable accuracy recovered by stopping.}

\bibliography{refs}

\end{document}
